\begin{table}[p]
    \centering
    \caption{\\ Gold clause exposure, investment, and payout by credit rating}
    \scriptsize
    \label{tab:credit_rating}

    \renewcommand{\arraystretch}{1.2}{
        \def\sym#1{\ifmmode^{#1}\else\(^{#1}\)\fi}
        \begin{tabular}{l*{2}{D{.}{.}{-1}}}
        \toprule
													                    & \multicolumn{1}{c}{Net investment} & \multicolumn{1}{c}{Dividend} \\
													                    & \multicolumn{1}{c}{(1)} & \multicolumn{1}{c}{(2)} \\
        \midrule
        \ensuremath{\text{1933} \times \tilde{d}} & -0.036\sym{***} & 0.025\sym{**} \\
         & (0.013) & (0.010) \\
        \ensuremath{\text{1934} \times \tilde{d}} & -0.047\sym{***} & 0.035\sym{***} \\
         & (0.008) & (0.010) \\
        \ensuremath{\text{1933} \times \text{Low rating}} & -0.038\sym{**} & 0.013\sym{***} \\
         & (0.016) & (0.005) \\
        \ensuremath{\text{1934} \times \text{Low rating}} & 0.029\sym{***} & 0.011 \\
         & (0.003) & (0.007) \\
        \ensuremath{\text{1933} \times \tilde{d} \times \text{Low rating}} & -0.073 & -0.025 \\
         & (0.046) & (0.016) \\
        \ensuremath{\text{1934} \times \tilde{d} \times \text{Low rating}} & -0.077\sym{***} & -0.037\sym{***} \\
         & (0.024) & (0.009) \\
        \midrule
        Firm FE & \multicolumn{1}{c}{Yes} & \multicolumn{1}{c}{Yes} \\
        Year FE & \multicolumn{1}{c}{Yes} & \multicolumn{1}{c}{Yes} \\
        \ensuremath{R^2} & 0.244 & 0.789 \\
        Observations & \multicolumn{1}{r}{7,074$\phantom{000}$} & \multicolumn{1}{r}{7,074$\phantom{000}$} \\
        \bottomrule
        \end{tabular}
        }

    \vspace*{3mm} \justifying \noindent
\scriptsize{\textit{Notes.} This table reports results from panel regressions testing whether debt overhang effects are more severe for firms closer to default. It regresses net investment (column 1) and cash dividends (column 2) on $Q$, gold clause exposure $\tilde{d}$ (as defined in equation (\ref{eq:tilde_d})), year $\times$ $\tilde{d}$ interactions, and triple interactions with a low credit rating indicator. Low rating equals one if a firm's bond rating is Ba or below in 1930 (the median rating among bond issuers). To ease interpretation, $\tilde{d}$ is normalized by subtracting its median value among firms with positive exposure. All regressions include firm and year fixed effects. Only 1933 and 1934 interaction terms are reported; see Internet Appendix Table \ref{tabapp:credit_rating} for the full list of coefficients. All variables are winsorized at the 0.5\% and 99.5\% levels within each year. Standard errors in parentheses are two-way clustered by firm and year. $^{*}p<0.10$, $^{**}p<0.05$, $^{***}p<0.01$.}
\end{table}